% ===================================================================
% Appendix to Section 2 (Setting): proofs of the anti-unification lemmas,
% lggs of infinite sets, and the (WS) lemma with its remarks.
% Source: research/theory/T1 §1.3 (Lemmas 1.2, 1.3; unchanged by verification);
%         T7 §1.3 ((WS) and its equivalent form, as revised after verification).
% ===================================================================
\providecommand{\vars}{\mathrm{vars}}
\providecommand{\Ttrust}{T_{\mathrm{trust}}}

\section{Proofs for Section~\ref{sec:setting}}\label{app:setting}

\Cref{thm:setting:lgg} is due to \citet{plotkin1970note} and
\citet{reynolds1970transformational}, and the rank lemma is folklore in their
lattice of schemas. All other statements of \Cref{sec:setting} are proved here.

\subsection{Proof of \texorpdfstring{\Cref{lem:setting:closure}}{the closure lemma} (closure facts)}

(a) Let $D$ be the set of judgments with an $R$-derivation from $B$. $D$ contains
$B$ (one-node derivations) and is closed under $R$: if $P\step j\in R$ and
$P\subseteq D$, graft derivations of the finitely many members of $P$ under a new
root $j$. So $\Cl_R(B)\subseteq D$. Conversely, if $X\supseteq B$ is $R$-closed,
induction on the height of derivations shows that every node of every derivation
from $B$ lies in $X$; so $D\subseteq\Cl_R(B)$. Extensivity and monotonicity in
$B$ and in $R$ hold because a derivation from $B$ using steps of $R$ is also one
from any $B'\supseteq B$ using steps of any $R'\supseteq R$. Idempotence: a
derivation from $\Cl_R(B)$ becomes a derivation from $B$ after grafting a
derivation from $B$ above each leaf. Finitarity: a derivation has finitely many
leaves, so it is a derivation from a finite $B_0\subseteq B$.

(b) A single step $P\step j\in R$ is a derivation of $j$ from $P$, so
$R\subseteq\Sound(R)$, and monotonicity in the rule set gives
$\Cl_R(B)\subseteq\Cl_{\Sound(R)}(B)$. Conversely, $\Cl_R(B)$ is closed under
$\Sound(R)$: if $P\subseteq\Cl_R(B)$ and $j\in\Cl_R(P)$, then by (a)
$j\in\Cl_R(\Cl_R(B))=\Cl_R(B)$; it contains $B$, so it contains
$\Cl_{\Sound(R)}(B)$. Finally, $\Sound(A)=\{P\step j:j\in\Cl_A(P)\}$ depends
only on $\Cl_A$, so $\Cl_A=\Cl_R$ gives $A\subseteq\Sound(A)=\Sound(R)$; and
$\Sound(R)$ itself has the closure $\Cl_R$. \qed

\subsection{Proof of \texorpdfstring{\Cref{lem:setting:rank}}{the rank lemma} (rank)}

Throughout, $|t|$ is the number of symbol occurrences in $t$, metavariable
occurrences included, and $\vars(t)$ is the set of distinct metavariables of $t$.
A \emph{renaming} of a set $X$ of metavariables is a substitution whose
restriction to $X$ is an injective map into metavariables. If $\tau=\sigma\theta$
with $\theta$ a renaming of $\vars(\sigma)$, then $\sigma=\tau\theta^{-1}$, so
$\sigma\equiv\tau$; conversely, $\sigma\equiv\tau$ means equality up to such a
renaming, and then $\mu(\sigma)=\mu(\tau)$.

Let $\tau=\sigma\theta$. Let $x_1,\dots,x_k$ be the distinct metavariables of
$\sigma$, and let $x_i$ occur $m_i\ge1$ times in $\sigma$.

\emph{Size.} Replacing each of the $m_i$ occurrences of $x_i$ (one symbol) by
$\theta x_i$ (with $|\theta x_i|$ symbols) gives
\[
  |\tau|=|\sigma|+\sum_{i=1}^k m_i\big(|\theta x_i|-1\big).
\]

\emph{Metavariables.} Every metavariable of $\tau$ occurs in some $\theta x_i$, so
\[
  |\vars(\tau)|\le\sum_{i=1}^k|\vars(\theta x_i)|,
\]
with equality iff the sets $\vars(\theta x_i)$ are pairwise disjoint.

\emph{Difference.} Since $|\vars(\sigma)|=k$,
\begin{equation}\label{eq:app:setting:rank}
  \mu(\tau)-\mu(\sigma)
  \;\ge\;\sum_{i=1}^k\Big[m_i\big(|\theta x_i|-1\big)+1-|\vars(\theta x_i)|\Big].
\end{equation}
We bound each summand.
\begin{itemize}
\item If $\theta x_i$ is a metavariable, then $|\theta x_i|=1=|\vars(\theta x_i)|$
  and the summand is $0$.
\item Otherwise $\theta x_i$ contains at least one function or constant symbol and
  at least one occurrence of each of its metavariables, so
  $|\theta x_i|\ge|\vars(\theta x_i)|+1$. Using $m_i\ge1$, the summand is at least
  $(|\theta x_i|-1)+1-|\vars(\theta x_i)|\ge1$.
\end{itemize}
Hence $\mu(\tau)\ge\mu(\sigma)$.

\emph{Equality.} If $\mu(\tau)=\mu(\sigma)$, then every summand in
\eqref{eq:app:setting:rank} is $0$ and the inequality there is tight. The first
forces every $\theta x_i$ to be a metavariable; the second forces the sets
$\vars(\theta x_i)=\{\theta x_i\}$ to be pairwise disjoint, i.e.\ the
$\theta x_i$ to be pairwise distinct. So $\theta$ restricted to $\vars(\sigma)$ is
a renaming. Conversely, if it is a renaming, then $|\tau|=|\sigma|$ and
$|\vars(\tau)|=|\vars(\sigma)|$, so $\mu(\tau)=\mu(\sigma)$.

\emph{Chains.} Suppose $\sigma_{l+1}\succ\sigma_l$ strictly, i.e.\
$\sigma_l=\sigma_{l+1}\theta_l$ and $\sigma_l\not\equiv\sigma_{l+1}$. Then
$\theta_l$ is not a renaming of $\vars(\sigma_{l+1})$ (otherwise the two would be
equivalent), so $\mu(\sigma_l)>\mu(\sigma_{l+1})$, i.e.\
$\mu(\sigma_l)\ge\mu(\sigma_{l+1})+1$. Summing along
$\tau=\sigma_0\prec\sigma_1\prec\dots\prec\sigma_m$ gives
$\mu(\tau)\ge\mu(\sigma_m)+m$, and $\mu(\sigma_m)\ge0$. \qed

\subsection{Proof of \texorpdfstring{\Cref{lem:setting:recover}}{the recovery lemma} (when the lgg recovers the schema)}

\emph{Idea.} Substituting $\theta x$ for the occurrences of $x$ adds more symbols
than distinct metavariables unless $\theta x$ is a metavariable, which is the rank
lemma. Anti-unifying the $t_j$ copies $\sigma$ and puts
$\mathcal A(\theta_1x,\dots,\theta_nx)$ at $x$; this is a metavariable iff (R)
holds, and two of them coincide iff (D) fails. In detail:

Let $\sigma$ have metavariables $x_1,\dots,x_v$, let $\theta_1,\dots,\theta_n$ be
ground substitutions ($n\ge1$), and let $t_j=\sigma\theta_j$. For each
metavariable $x$ of $\sigma$ write $c_x=(\theta_1x,\dots,\theta_nx)$ for its
\emph{column}, and let $\eta$ be the substitution $\eta(x):=\mathcal A(c_x)$.

\emph{Claim.} For every subterm $\sigma'$ of $\sigma$,
$\mathcal A(\sigma'\theta_1,\dots,\sigma'\theta_n)=\sigma'\eta$.

By induction on $\sigma'$. If $\sigma'=x$ is a metavariable, both sides are
$\mathcal A(c_x)$. If $\sigma'=f(\sigma'_1,\dots,\sigma'_r)$ (with $r=0$ for a
constant), then every $\sigma'\theta_j=f(\sigma'_1\theta_j,\dots,\sigma'_r\theta_j)$
has root $f$, so by definition $\mathcal A$ outputs $f$ applied to
$\mathcal A$ of the argument columns $(\sigma'_k\theta_1,\dots,\sigma'_k\theta_n)$,
which by the induction hypothesis are $\sigma'_k\eta$. Hence the output is
$f(\sigma'_1\eta,\dots,\sigma'_r\eta)=\sigma'\eta$.

Taking $\sigma'=\sigma$ and using \Cref{thm:setting:lgg},
\[
  \lgg(t_1,\dots,t_n)=\mathcal A(t_1,\dots,t_n)=\sigma\eta\preceq\sigma .
\]

\emph{When $\eta$ is a renaming.} By the definition of $\mathcal A$:
\begin{itemize}
\item $\eta(x)=\mathcal A(c_x)$ is a metavariable iff the entries of $c_x$ do not
  all have the same root, i.e.\ iff (R) holds at $x$. If (R) fails at $x$,
  $\eta(x)$ is a compound term or a constant.
\item If (R) holds at $x$ and at $y$, then $\eta(x)=z_{c_x}$ and $\eta(y)=z_{c_y}$,
  and these coincide iff $c_x=c_y$, i.e.\ iff (D) fails for the pair $x,y$.
\end{itemize}
So $\eta$ restricted to $\vars(\sigma)$ is a renaming iff (R) holds for every $x$
and (D) holds for every pair $x\neq y$.

\emph{Conclusion.} If (R) and (D) hold, $\eta$ is a renaming and
$\lgg(t_1,\dots,t_n)=\sigma\eta\equiv\sigma$. If either fails, $\eta$ is not a
renaming of $\vars(\sigma)$, so by \Cref{lem:setting:rank}
$\mu(\sigma\eta)>\mu(\sigma)$; since equivalent schemas have equal rank,
$\lgg(t_1,\dots,t_n)\not\equiv\sigma$, and it is strictly more specific. \qed

\subsection{Least general generalizations of infinite sets}\label{app:setting:infinite}

Every nonempty set $P$ of ground terms, finite or not, has a least general
generalization, unique up to renaming. To see this, choose a nonempty finite
$Q\subseteq P$ for which $\mu(\lgg(Q))$ is minimal ($\mu$ takes values in $\N$).
Let $p\in P$. Since $Q\subseteq Q\cup\{p\}\subseteq\inst(\lgg(Q\cup\{p\}))$,
\Cref{thm:setting:lgg} gives $\lgg(Q)\preceq\lgg(Q\cup\{p\})$, so by
\Cref{lem:setting:rank} $\mu(\lgg(Q\cup\{p\}))\le\mu(\lgg(Q))$. By minimality the
two ranks are equal, so by the equality case of \Cref{lem:setting:rank}
$\lgg(Q\cup\{p\})\equiv\lgg(Q)$ and $p\in\inst(\lgg(Q))$. Hence
$P\subseteq\inst(\lgg(Q))$. If $P\subseteq\inst(\tau)$, then
$Q\subseteq\inst(\tau)$ and $\lgg(Q)\preceq\tau$. So $\lgg(Q)$ is a least general
generalization of $P$; any two are $\preceq$ each other, hence equivalent. \qed

\subsection{Proof of \texorpdfstring{\Cref{lem:setting:guard}}{the guarded closure lemma} (guarded closure)}

$P\subseteq\inst(\lgg(P))$ by \Cref{thm:setting:lgg}, and every $\phi\in\Psi_P$
holds on $P$, so $P\subseteq\inst(\lgg(P),\Psi_P)$. Let
$P\subseteq\inst(\sigma,\Psi)$. Then $P\subseteq\inst(\sigma)$, so
$\sigma\succeq\lgg(P)$ and $\inst(\lgg(P))\subseteq\inst(\sigma)$. Every
$\phi\in\Psi$ holds on all of $P$, so $\Psi\subseteq\Psi_P$ and
$\bigcap_{\phi\in\Psi_P}\phi^{-1}(1)\subseteq\bigcap_{\phi\in\Psi}\phi^{-1}(1)$.
Intersecting, $\inst(\lgg(P),\Psi_P)\subseteq\inst(\sigma,\Psi)$. \qed

\subsection{Sound evaluation and coherence}\label{app:setting:WS}

\begin{lemma}[(WS) and coherence; trivial]\src{T7 \S1.3}\label{lem:setting:WS}
Under (WS) every designated position is in bounds for $\Rstar\cup\Ttrust$, hence
for $\Rstar$. Conversely, if $J_W=\emptyset$ and every designated position is in
bounds for $\Rstar\cup\Ttrust$ (for instance, in bounds for $\Rstar$ with
$\Ttrust\subseteq\Sound(\Rstar)$), then (WS) holds.
\end{lemma}

\emph{Remarks.} (WS) does not require $\Ttrust\subseteq\Sound(\Rstar)$: in
arithmetic the target may consist of non-logical axioms with first-order logic
trusted (\Cref{sec:twotier}), and the relevant target closure is then
$\Cl_{\Rstar\cup\Ttrust}$. In formal mathematics $V$ is truth in the intended
structure $\mathfrak M$: (WS) holds when designated positions are true of
$\mathfrak M$, $W$ is truth in $\mathfrak M$ on its fragment, and target and
trusted steps are $\mathfrak M$-sound.

\begin{proof}
Write $R^+=\Rstar\cup\Ttrust$. Let $\pos{\Gamma}{\Delta}$ be designated and $V$ as
in (WS). $V$ makes every member of $\Gamma$ true, and no step of $R^+$ has true
premises and a false conclusion, so the set of $V$-true judgments is an
$R^+$-closed superset of $\Gamma$ and contains $\Cl_{R^+}(\Gamma)$. Since $V$
makes $\bot$ and every member of $\Delta$ false, $\bot\notin\Cl_{R^+}(\Gamma)$
and $\Delta\cap\Cl_{R^+}(\Gamma)=\emptyset$. As $\Cl_{\Rstar}(\Gamma)\subseteq
\Cl_{R^+}(\Gamma)$ (\Cref{lem:setting:closure}(a), monotonicity in the rule set),
the position is also in bounds for $\Rstar$.

Conversely, let $J_W=\emptyset$ and let $\pos{\Gamma}{\Delta}$ be a designated
position that is in bounds for $R^+$. Then $e$ is well defined, because
$\Gamma\subseteq\Cl_{R^+}(\Gamma)$ is disjoint from $\Delta\cup\{\bot\}$. Let $V$
be the indicator function of $\Cl_{R^+}(\Gamma)$. It extends $e$: it is $1$ on
$\Gamma$ and $0$ on $\Delta\cup\{\bot\}$. The set $\Cl_{R^+}(\Gamma)$ is closed
under $R^+$, so no step of $R^+$ has $V$-true premises and a $V$-false conclusion,
and (WS) holds. For the parenthetical instance: if $\Ttrust\subseteq\Sound(\Rstar)$,
then $\Rstar\subseteq R^+\subseteq\Sound(\Rstar)$, so by
\Cref{lem:setting:closure}(a),(b) $\Cl_{R^+}=\Cl_{\Rstar}$, and being in bounds
for $\Rstar$ is the same as being in bounds for $R^+$.
\end{proof}
